%===============================================================================
% CHAPTER 1: INTRODUCTION
%===============================================================================
\chapter{Introduction}
\label{ch:introduction}

\section{Background and Motivation}
\label{sec:background}

Democratic processes rely fundamentally on the integrity of voting systems. Whether for
national elections, student councils, or organizational governance, the ability to
cast a vote and trust that it will be counted correctly is paramount. Traditional
paper-based voting, while providing tangible audit trails, suffers from high costs,
logistical complexity, slow tallying, and accessibility barriers for remote voters.

Electronic voting (e-voting) emerged to address these limitations, offering speed,
convenience, and potential cost savings. However, early e-voting systems—particularly
Direct Recording Electronic (DRE) machines and centralized internet voting
platforms—introduced new vulnerabilities:
\begin{itemize}
    \item \textbf{Centralized Trust:} A single database administrator can alter results.
    \item \textbf{Opacity:} Voters cannot independently verify their vote was recorded.
    \item \textbf{Single Points of Failure:} Server compromise endangers the entire election.
    \item \textbf{Lack of Transparency:} Closed-source code prevents public audit.
    \item \textbf{Coercion Vulnerability:} No receipt-freeness guarantees.
\end{itemize}

High-profile incidents—from the 2000 US presidential election controversies to
vulnerabilities discovered in Estonia's internet voting system—have eroded public
confidence. The core problem remains: \textit{how to achieve the efficiency of digital
systems without sacrificing the verifiability and trust of physical ballots}.

\section{Problem Statement}
\label{sec:problem}

Existing blockchain-based voting proposals often fall into two extremes:
\begin{enumerate}
    \item \textbf{Smart-contract platforms (Ethereum, etc.):} High gas costs (\$5--50/vote),
    low throughput (15--30 TPS), and complex user experience (wallet management).
    \item \textbf{Permissioned ledgers (Hyperledger Fabric):} Require consortium governance,
    defeating the purpose of public verifiability.
\end{enumerate}

A practical e-voting system for real-world organizations (universities, associations,
cooperatives) needs:
\begin{itemize}
    \item \textbf{Cryptographic immutability:} Votes cannot be altered post-submission.
    \item \textbf{Public auditability:} Anyone can verify the tally independently.
    \item \textbf{Low cost:} Sub-cent per vote to enable large-scale elections.
    \item \textbf{High throughput:} Support thousands of concurrent voters.
    \item \textbf{Familiar UX:} Web-based, no cryptocurrency wallet required.
    \item \textbf{Regulatory alignment:} Identity verification compatible with KYC/AML.
\end{itemize}

\section{Objectives}
\label{sec:objectives}

\subsection{Primary Objective}
Design and implement a \textbf{hybrid cloud and distributed ledger e-voting system}
that achieves tamper-proof vote recording, transparent tallying, and practical
usability for organizational elections.

\subsection{Secondary Objectives}
\begin{enumerate}
    \item Implement vote immutability via Hedera Consensus Service (HCS) topics.
    \item Enforce one-person-one-vote through atomic database-ledger coordination.
    \item Provide real-time transparent results via Hedera mirror nodes.
    \item Achieve sub-5-second finality with predictable \$0.0001 per-vote cost.
    \item Build an informed-voting interface with candidate manifestos and evidence.
    \item Develop administrative tooling for election lifecycle management.
\end{enumerate}

\section{Scope and Limitations}
\label{sec:scope}

\subsection{In Scope}
\begin{itemize}
    \item Student council / organizational elections (not national elections)
    \item Web-based React single-page application
    \item Hedera Testnet integration (production-ready architecture, testnet deployment)
    \item Role-based access: Administrators and Voters
    \item Manifesto dashboard with campus issues and evidence documents
    \item Docker-based deployment for reproducibility
\end{itemize}

\subsection{Out of Scope}
\begin{itemize}
    \item National-level election regulatory compliance
    \item Biometric or hardware-based authentication
    \item Zero-knowledge proofs for vote privacy (pseudonymous only)
    \item Hedera Mainnet deployment (requires governance review)
    \item Native mobile applications (responsive web only)
    \item Vote delegation / liquid democracy
    \item Post-quantum cryptographic migration
\end{itemize}

\section{Report Organization}
\label{sec:organization}

This report is structured as follows:
\begin{itemize}
    \item \textbf{Chapter 2} reviews related work in e-voting, blockchain voting,
    and hybrid architectures.
    \item \textbf{Chapter 3} presents the system architecture, component design,
    data models, and security model.
    \item \textbf{Chapter 4} details the implementation of backend services,
    frontend components, and blockchain integration.
    \item \textbf{Chapter 5} reports functional testing, performance metrics,
    security evaluation, and usability results.
    \item \textbf{Chapter 6} discusses achievements, challenges, limitations,
    and future work.
    \item \textbf{Chapter 7} concludes with contributions and impact.
    \item \textbf{Appendices} provide source code structure, API documentation,
    deployment guide, user manuals, and test cases.
\end{itemize}